% Resumo: de 150 a 250 palavras (template). Contexto, motivação, objetivo, justificativa,
% metodologia e resultados alcançados.
\textit{Redes de casas e de pequenos negócios reúnem dispositivos da Internet das Coisas (IoT) e não têm especialistas em segurança. Os sistemas de detecção com aprendizado de máquina param no alerta, e as plataformas de detecção e resposta de rede (NDR, Network Detection and Response) são fechadas. Este trabalho apresenta o \mbox{SentryIoT}, um sistema de NDR aberto que detecta ataques com um classificador Random Forest e responde a eles com três agentes de modelos de linguagem de grande escala (LLM, Large Language Model), num modelo local. Um servidor Model Context Protocol (MCP) aplica a política de risco e exige aprovação humana nas ações de risco alto. Treinado com dados regerados dos pcaps do CICIoT2023, o classificador reconheceu mais de 98\% das janelas de negação de serviço e de Mirai no trecho de teste de cada captura; varredura, ataques Web e força bruta ficaram entre 34,4\% e 55,0\%. Em uma prova de conceito, com pcaps reproduzidos e ações simuladas, os agentes trataram 28 incidentes abertos pelo classificador, 19 deles verdadeiros, cada um três vezes. Com o limiar de confiança, a triagem reconheceu os 19, e os agentes aplicaram uma ação aceitável no alvo certo em 16 deles, com mediana de 60,86~s até a remediação. Nenhuma ação atingiu o gateway nem os serviços críticos da rede. Nesta prova de conceito, um modelo local levou os incidentes do alerta à resposta, inclusive decidindo com rótulos no lugar dos endereços da rede, com as garantias de segurança aplicadas fora do modelo.}

\vspace{\onelineskip}

\noindent
\textbf{Palavras-chave}: Internet das Coisas; sistemas multiagentes; detecção e resposta de rede; Random Forest; modelos de linguagem de grande escala.
